\section{Project overview}
\label{sec:overview}

\subsection{What the workspace contains}
The repository has a reusable library and in-tree applications with one-way
dependency flow.
\textbf{Illumo} is a static C++ library containing the reusable host,
application runner, platform layer, foundation, services, rendering, generic
console/UI, persistent scene hierarchy, assets, and Debug module.
\textbf{IllumoGame} is the interactive cellular-automata simulator that consumes
Illumo through supported \texttt{<Illumo/...>} headers. It owns the canvas,
rulesets, simulation and persistence policy, product commands, CA defaults, and
a declarative application definition. \textbf{IllEd} is the SceneGraph world
editor that writes \texttt{.ilsc} documents so later Illumo applications can
be authored rather than only coded. \textbf{IllMeshViewer} provides interactive
mesh and scene inspection. These products do not contain process entry,
native dialog implementation, logging lifetime, system parser, or frame loop.

\textbf{Illumo::Content} (D-E18) is an optional library above the engine that
every product links. It owns the one scene format, \texttt{.ilsc} format 2
(D-E19), and \symbolref{SceneInstance}, the live loader that builds a
SceneGraph and its attachments from a document. IllEd edits scenes through it,
IllMeshViewer opens them, and IllumoGame's \texttt{render3dTest} diagnostic is
one. Every program is a package described by \texttt{illumo.json} (D-E22),
either a loose directory or an \texttt{.ilpk} archive (D-E20). The runtime
mounts packages on one virtual file tree (D-E21): the launched application at
\texttt{/app}, engine assets at \texttt{/engine}, every other installed or
\texttt{--mount}ed package at \texttt{/packages/<id>}, and an optional writable
\texttt{/project}. Mods and expansions are additive: a content package adds
files under its own mount and may overlay files onto \texttt{/app}, but cannot
delete a base file. Guests reach the tree through file protocol v2 (D-E23) and
a fetch-on-demand asset cache (D-E24).

IllEd (D-E25) has patch undo/redo, a multi-node selection with box select,
translate/rotate/scale gizmos in local or world space with snapping, a typed
inspector with text entry, clipboard copy/paste of subtrees, a hierarchy with
reordering and folding, mesh/sprite/light/camera nodes, an environment editor,
and an asset browser over the file tree. With \texttt{--project <dir>} it
imports files into the project, saves scenes there, and packs the project into
an \texttt{.ilpk}. Its Hierarchy, Assets, Tools and Inspector panels, like
IllMeshViewer's Info and Display panels, dock in columns or pop out into
their own windows, in a plain tool look (D-UI7).

IllumoGame supports multiple rule sets (Game of Life, Brian's Brain,
Wireworld, \ldots), camera navigation, edit/run modes, save/load, and a compact
in-app command/debug surface. Rendering is OpenGL via GLFW/GLEW on the supported
Windows path. Linux sources are repaired for Ubuntu 24.04 x86\_64
X11/XWayland; they are not promoted to supported until native smoke exists.
macOS is not targeted; no platform scaffold is retained.

\begin{inferencebox}
Illumo is intended to be a reusable runtime and rendering foundation for
future projects, with IllumoGame as its first in-tree consumer. D-E6 records
the explicit library/product boundary and D-E8 adds a bounded persistent scene
hierarchy. This direction does not imply a full ECS, render graph, stable DLL
ABI, or speculative multi-backend framework.
\end{inferencebox}

\subsection{High-level system diagram}

\begin{center}
\begin{tikzpicture}[
  font=\small,
  box/.style={draw, rounded corners, align=center, inner sep=6pt, minimum height=1.1cm},
  arr/.style={-{Latex}, thick}
]
  \node[box, fill=orange!10] (game) {IllumoGame / IllEd\\IllMeshViewer\\application definitions};
  \node[box, fill=green!10, right=2.0cm of game] (lib) {Illumo static library\\Application + Platform + Engine\\Scene + Foundation + Services + Rendering\\(+ Illumo::Content)};
  \node[box, fill=purple!8, below=1.0cm of lib] (gl) {private OpenGL backend\\Illumo assets + licenses};
  \node[box, fill=gray!10, below=1.0cm of game] (tests) {Four test runners\\Illumo::TestSupport};

  \draw[arr] (game) -- node[above]{\scriptsize \texttt{Illumo::Illumo}} (lib);
  \draw[arr] (lib) -- (gl);
  \draw[arr] (tests) -- (game);
  \draw[arr] (tests) -- (lib);
\end{tikzpicture}
\end{center}

\subsection{Runtime one-liner}
Illumo platform entry $\rightarrow$ \symbolref{CreateIllumoApplication}
$\rightarrow$ engine-owned \symbolref{RunIllumoApplication} applies supplied
product defaults and metadata-driven CLI $\rightarrow$ construct and fallibly
initialize \symbolref{Illumo} $\rightarrow$ register the supplied required product
module and engine-owned optional Debug module $\rightarrow$ \texttt{std::chrono} loop
\{\,\symbolref{update}, \symbolref{render}\,\} until window close.

\subsection{Non-goals (current)}
\begin{itemize}
  \item No install/export package, shared DLL, or stable binary ABI.
  \item No general-purpose AAA renderer, full ECS, or speculative framework.
  \item No SceneGraph self-serialization (scenes are written by
        Illumo::Content's codec through the public graph API), retained UI,
        culling hierarchy, or update framework until a concrete consumer
        defines those contracts.
  \item No material system, animation, scripting, physics, prefabs, or nested
        scene instances in the scene format yet.
  \item No multiplayer or networking.
  \item No Linux support claim until native validation exists; macOS is not targeted.
\end{itemize}
